%; whizzy chapter
% -initex iniptex -latex platex -format platex -bibtex jbibtex -fmt fmt
% 以上 whizzytex を使用する場合の設定。

%     Tokyo Debian Meeting resources
%     Copyright (C) 2011 Junichi Uekawa
%     Copyright (C) 2011 Nobuhiro Iwamatsu

%     This program is free software; you can redistribute it and/or modify
%     it under the terms of the GNU General Public License as published by
%     the Free Software Foundation; either version 2 of the License, or
%     (at your option) any later version.

%     This program is distributed in the hope that it will be useful,
%     but WITHOUT ANY WARRANTY; without even the implied warranty of
%     MERCHANTABILITY or FITNESS FOR A PARTICULAR PURPOSE.  See the
%     GNU General Public License for more details.

%     You should have received a copy of the GNU General Public License
%     along with this program; if not, write to the Free Software
%     Foundation, Inc., 51 Franklin St, Fifth Floor, Boston, MA  02110-1301 USA

%  preview (shell-command (concat "evince " (replace-regexp-in-string "tex$" "pdf"(buffer-file-name)) "&"))
% 画像ファイルを処理するためにはebbを利用してboundingboxを作成。
%(shell-command "cd image201109; ebb *.png")

%%ここからヘッダ開始。

\documentclass[mingoth,a4paper]{jsarticle}
\usepackage{monthlyreport}

% 日付を定義する、毎月変わります。
\newcommand{\debmtgyear}{2011}
\newcommand{\debmtgmonth}{9}
\newcommand{\debmtgdate}{17}
% (+ (* (- 2011 2005) 12) 9 -1) started from zero
\newcommand{\debmtgnumber}{80}

\begin{document}

\begin{titlepage}
\thispagestyle{empty}
% タイトルページ:編集必要な部分は最初のマクロに飛ばすこと

\vspace*{-2cm}
第\debmtgnumber{}回 東京エリア Debian 勉強会資料\\
\hspace*{-2cm}
\includegraphics[width=210mm]{image201003/debsen.eps}\\
\hfill{}\debmtgyear{}年\debmtgmonth{}月\debmtgdate{}日

% ここはアップデートすること
\rotatebox{10}{\fontsize{32}{32} {\gt 特集1: XXXX}}

\rotatebox{10}{\fontsize{32}{32} {\gt 特集2: YYYY}}

\vspace*{-2cm}
\hfill{}\includegraphics[height=6cm]{image200502/openlogo-nd.eps}
\end{titlepage}

\dancersection{Introduction}{上川 純一}

\begin{multicols}{2}
 

 今月のDebian勉強会へようこそ。これからDebianの世界にあしを踏み入れると
 いう方も、すでにどっぷりとつかっているという方も、月に一回Debianについ
 て語りませんか？

 Debian勉強会の目的は下記です。

 \begin{itemize}
 \item \underline{Debian Developer} (開発者)の育成。
 \item 日本語での「\underline{開発に関する情報}」を整理してまとめ、アップデートする。
 \item \underline{場}の提供。
 \begin{itemize}
  \item 普段ばらばらな場所にいる人々が face-to-face で出会える場を提供
	する。
  \item Debian のためになることを語る場を提供する。
  \item Debianについて語る場を提供する。
 \end{itemize}
 \end{itemize}		

 Debianの勉強会ということで究極的には参加者全員がDebian Packageをがりがり
 と作るスーパーハッカーになった姿を妄想しています。情報の共有・活用を通し
 て Debianの今後の能動的な展開への土台として、「場」としての空間を提供す
 るのが目的です。

\end{multicols}

\newpage

\begin{minipage}[b]{0.2\hsize}
 \definecolor{titleback}{gray}{0.9}
 \colorbox{titleback}{\rotatebox{90}{\fontsize{80}{80} {\gt デビアン勉強会} }}
\end{minipage}
\begin{minipage}[b]{0.8\hsize}
\hrule
\vspace{2mm}
\hrule
\begin{multicols}{2}
\tableofcontents
\end{multicols}
\vspace{2mm}
\hrule
\end{minipage}

\dancersection{事前課題}{岩松 信洋}

今回の事前課題は以下です:
\begin{enumerate}
 \item Debian使いとしてウェブサービスに期待することは何ですか？
\end{enumerate}
この課題に対して提出いただいた内容は以下です。
\begin{multicols}{2}
{\small
% \begin{prework}{ ����(yy\_y\_ja\_jp) }
���ޤ���������ǤϤʤ��Ǥ���... �긵�Υ�åץȥåפǤ� /boot �ˤ� ext2
 ����¾�ˤ� ext3���Ƕ�ȤäƤ���ǥ����ȥåפǤ� /boot �ˤ� ext3 ����¾
 �ˤ� LVM ��� ext4 ��ȤäƤ��ޤ���
\end{prework}

\begin{prework}{ �����ϥ� }
�ǥե���Ȥ�ext3����
�����ǥե���ȤΤޤޡ�
(NTFS���VM���᡼�����ext3�⤢�뤱��)
\end{prework}

\begin{prework}{ yos.takahashi }
ext3/4���˻ȤäƤޤ���ext3�Υǡ�������١����ˤĤ�������Linux2011ǯ1���˼�ɮ���ޤ�����
\end{prework}

\begin{prework}{ MATOHARA }
����inode �ϳ���������Ƥ���inode ��ưŪ�˳�����Ƥ���XFS �����򤹤뤳
 �Ȥ�¿���Ǥ���NILFS �Ͼ�����Ƥߤ��ΤǤ�����mount ���˰ʲ��Τ褦�ʥ��
 ���������ФƤޤ��ݤ��ʤȻפ��ޤ�����
\begin{commandline}
$ sudo mount /dev/sdb1 /mnt
mount.nilfs2: WARNING! - The NILFS on-disk format may change at any time.
mount.nilfs2: WARNING! - Do not place critical data on a NILFS filesystem. 
\end{commandline}
����¾NotePC �Ǥ�dm-crypt �ξ�˥ե����륷���ƥ���֤��ưŹ沽�����ꡢ
 eCryptfs �ǰŹ沽�����ꤷ�Ƥ��ޤ����񤭹��߻���CPU �򤫤ʤ���񤷤ޤ��ġ�
\end{prework}

\begin{prework}{ ��ޤ� }
��ext2-$>$reiserfs-$>$jfs-$>$xfs-$>$reiserfs-$>$ext3�ȻȤäƤ��ޤ�����

����:
\begin{itemize}
 \item reiserfs: ���ե������¿���ե�����Υ������������ӥ��Ӥ��Ƥ��ɤ���
       �����������դȤ����ⵤ�������θ�μ�žȬ�ݥ����ɤˡ�����
 \item jfs: ����v1.0��̾��ä�IBM���ꥨ�ʥ��������ˤ�xfs��ƨ��
 \item xfs: fsck==true�˴�ư����������ǯ�Ȥä���ΤΥޥ�����Ĵ����0byte
       �ե��������������Ѥ���줺ƨ����
\end{itemize}
������reiser4��Ķ���Ԥ��뤦���ˤ��줬�����ʤäơ����ext3�˸��경��htree�����ä������⤦Ŵ�Ĥʤ鲿�Ǥ⤤���Ǥ��������Ȥ����Ĥ�nilfs�ʤɤ˼��Ф��Ƥ��ޤ���ext3�������noatime���٤Ǥ����������aufs���碌�Ƽ�ʬ���Ȥ��Ȥ� *strap �Ķ��򥯥����˥󥰤��Ƥ��ޤ����¸����ƥ��Ȥ������Ǥ�����Ǥ���USB�����ư�Ǥ�ͭ�ѡ�

���LVM�ǤϤʤ�MD��Ȥäƾ�Ĺ���ܥХå����åפ򤷤Ƥ��ޤ���
 MD(sda,sdb,sdc)�ǹ��ۤ����̾��MD(����)�Dz�ư���Хå����åפλ���
 attach/detach�򤹤롣�֥��å���٥�ʤΤǥꥫ�Х��FSǤ���Ǥ��������̤�
 �������¾����ˡ���ʤ�������
\end{prework}

\begin{prework}{ henrich }
�ȤäƤ���֤�NTFS��Ĺ���󤸤�ʤ��Ǥ����͡����졣
�����Ρ����Ѥο������ǥ�������ext4�ǥե����ޥåȤ��ޤ����������㤤��Ƚ��ʤ��Ȥ����������Ƥ��ޤ���
\end{prework}

\begin{prework}{ emasaka }
�Ĥ뤷��FS��ȤäƤޤ�
\end{prework}

\begin{prework}{ �ܾ� }
ext3����Ѥ��Ƥ��ޤ����ä��Ѥ�ä����ȤϤ��Ƥ��ޤ���
FS����ʤ��Ǥ������Ƕ�Lenny��2TB��HDD��Ȥä���parted�äƤλȤ��ƶä��ޤ�����
\end{prework}

\begin{prework}{ ����@������ }
���������Ū�˳��Ѥ��Ƥ���ե����륷���ƥ��ReiserFS�Ǥ���
�Ż��ǻȤäƤ�Ķ���ext3�Ǥ�����ext3���Ÿ������ǥ��㡼�ʥ뤬
����Ʋ��Ǥ������ηи��ʸŤ������ͥ�Ǥ���...�ˤ����ꡢ
���ޤ꿮�Ѥ��Ƥ��ޤ���
�����ReiserFS�Ķ��ǤϤ���ޤǤνꤤ���ʤ��Ÿ����ڤä��ꡢ
�Ƥ�HDD�����줫�����ꤷ�Ƥ��ﳲ����ä��и���̵���Τǡ�
��³Ū�˻ȤäƤ��ޤ���
��ǯ����ReiserFS�Υᥤ��ȯ��(Hans Reiser)�����ᤵ��Ƥ��ޤ������ƥʥ󥹤��ۤ��Ƥ��ޤ�����
�����������θ��ReiserFS��¾�γ�ȯ�Ԥˤ���³�����ݼ餵��Ƥ���Τǡ��¿����ޤ�����
\end{prework}

\begin{prework}{ nozzy123nozzy }
\begin{enumerate}
 \item LVM�ˤĤ��Ƥϡ�CentOS5.5��Ƴ�������Τ��Τޤޤ����Ѥ��Ƥޤ�����
       ���������ƥब��äƤ���Volume̾�ϥǥե���Ȥ�����ȡʼºݤˤ�
       kickstart�ˤơˤ��ѹ����ƻȤäƤޤ����ʾ㳲���Υ���١����˺��뤿
       ���
\item ext3�ˤĤ��Ƥϡ�debian-sid�����Τޤ޻��ꤷ�Ƥ����Τ򤽤Τޤ޻Ȥ�
      �Ƥ����ꤷ�ޤ���relatime, noatime ���餤�Ͼ����ɲä��Ƥߤ����ʡ���
      �ϻפäƤޤ���
\end{enumerate}
\end{prework}

\begin{prework}{ �ޤ��������ؤ� }
\begin{itemize}
 \item Debian�Ǥ��ä˶Ťä����Ȥ�����ext3��ȤäƤޤ������ۥޥ����qcow2
       ���᡼���ǥ��������ѻ��ʳ��ϡ�LVM�ϻȤäƤޤ���
 \item �����ǰ����ü�ʤΤϡ��������DHCP�������Ѥ�Armadillo-J�ǻȤäƤ�
       ��JFFS�Ǥ����ǥե���ȤΥե����०�����Ǥϥ�֡��Ȥ�����������
       �ƽ��������Ƥ��ޤ��Τǡ�RAM�ΰ�˽񤭤��ߡ��Ÿ��ڤäƤ�ä��ʤ�
       ���������Ǥ���Debian��udhcp�Υ������ѥå���������ӥ�ɤ��ƻȤäƤޤ���
       \footnote{\url{http://d.hatena.ne.jp/mkouhei/20080601/1212330630}}
 \item ��Debian���ߤǡ���ʬ����ǰ��֥ۥåȤʤΤ�palm webOS�Ǥ��������Ubuntu��
       �������ޥ���������Τ餷���ΤǤ�����/etc/mtab�򸫤��35�Ԥ⤢�ꡢ
       ���ʤ����֤ʹ����Ǥ��͡�
\end{itemize}
\end{prework}
}
\end{multicols}

\dancersection{最近のDebian関連のミーティング報告}{岩松 信洋}
\subsection{東京エリアDebian勉強会79回目報告}

% (query-replace-regexp "<.*?>" "")
% (query-replace-regexp "^[	 ]\+" "")



\printindex

\cleartooddpage

\vspace*{15cm}
\hrule
\vspace{2mm}
\includegraphics[width=2cm]{image200502/openlogo-nd.eps}
\noindent \Large \bf Debian 勉強会資料\\
\noindent \normalfont \debmtgyear{}年\debmtgmonth{}月\debmtgdate{}日 \hspace{5mm}  初版第1刷発行\\
\noindent \normalfont 東京エリア Debian 勉強会 （編集・印刷・発行）\\
\hrule

\end{document}
